\begin{titlepage}
    \vbox{ }

    \vbox{ }

    \begin{center}
        % Over-titre
        \includegraphics[width=0.40\textwidth]{img/logos/logo.png}\\[1cm]
        \textsc{\Large Rapport}\\[0.2cm]
        \textsc{\Large Projet de programmation - Semestre 2}\\[0.6cm]

        % \vbox{ }
        % Titre
        \noindent\makebox[\linewidth]{\rule{.7\paperwidth}{.6pt}}\\[0.7cm]
        { \huge \bfseries Nyctalopia}\\[0.25cm]
        \noindent\makebox[\linewidth]{\rule{.7\paperwidth}{.6pt}}\\[0.7cm]
        \large{\bfseries Créé par :}\\[0.1cm]
        \large{Iñigo Aldaraborda Hoyos - Chef de projet}\\[0.1cm]
        \large{Hugo Meleiro}\\[0.1cm]
        \large{Guilhem Cros}\\[0.1cm]
        \large{Marin Godefroy}\\[0.1cm]
        \vfill
        % Formatter
        \large
        % Lieu
        EPITA Toulouse

            % Date
            {\large Mardi 26 Avril 2022}
    \end{center}
\end{titlepage}
